\fboxsep0.5pt %make colorboxes smaller with this
\begin{frame}[fragile]{Stellingen en Bewijzen (Nummering)}
    Vaak heb je dat je zowel lemma's als stellingen in je bestand hebt. \pause

    Standaard hebben deze verschillende nummering zodat je 
    Stelling 1 en Lemma 1 krijgt. \pause 

    Meestal wil je dat de nummering doorloopt, dus je wil Lemma 1, Stelling 2, etc. 
    Om dit te doen gebruiken we:
    \begin{minted}{tex}
        \newtheorem{stl}{Stelling}
        \newtheorem{lem}|\colorbox{cyan!40!white}{[stl]}|{Lemma}
    \end{minted}
    Nu loopt de teller van `Lemma' mee met die van Stelling. 
\end{frame}

\begin{frame}[fragile]{Stellingen en Bewijzen (Nummering)}
    Je kunt ook je stellingen meenummeren met je hoofdstukken. 
    \pause
    Je doet dit met:
    \begin{minted}{tex}
        \newtheorem{stelling}{Stelling}|\colorbox{cyan!40!white}{[section]}|
        \newtheorem{lemma}[stelling]{Lemma}
    \end{minted}
    Kijk goed waar de \texttt{[...]} staat!
    \pause
    Nu krijg je in \texttt{section 1} bijvoorbeeld Stelling 1.1, Lemma 1.2, etc. 
\end{frame}